\chapter{The Expanded Welch--Satterthwaite Test for MI}
\label{chap:expanded}

Expanded Welch keeps the estimated MI difference, standard error and test
statistic from Chapter~\ref{chap:baselines}. It changes the distribution used
to judge that statistic from the standard normal distribution to a Student
\(t\) distribution. To choose its degrees of freedom, this chapter calculates
how much the estimated MI variance changes between samples.

\section[Choosing Degrees of Freedom for Expanded Welch]{Choosing Degrees of
Freedom for\\Expanded Welch}

The Welch--Satterthwaite equation combines the two estimated MI variance
contributions
\citep{satterthwaite1946,welch1947}:
\begin{equation}
  \nu
  =
  \frac{
  \left\{\Vhat(P)/n_P+\Vhat(Q)/n_Q\right\}^2
  }{
  \left\{\Vhat(P)/n_P\right\}^2/\nu_P
  +\left\{\Vhat(Q)/n_Q\right\}^2/\nu_Q
  }.
  \label{eq:ws-general}
\end{equation}
The variance contributions \(\Vhat(P)/n_P\) and \(\Vhat(Q)/n_Q\) already
appear in the shared standard error. The new quantities are the component
degrees of freedom \(\nu_P\) and \(\nu_Q\). They describe how much each
estimated MI variance varies between samples relative to its mean.

The derivation has four steps. First, we relate the degrees of freedom to the
average and across-sample variation of \(\Vhat(P)\). Second, we calculate how \(V(P)\)
changes when a small amount of probability moves into one cell. Third, that
change tells us how variable \(\Vhat(P)\) is. Finally, we combine the
degrees of freedom from \(P\) and \(Q\) using Equation~\eqref{eq:ws-general}.

\medskip
\noindent\textbf{Moment matching for one component.}

Satterthwaite approximates the positive variance estimate \(\Vhat(P)\) by
the scaled chi-squared variable
\begin{equation}
  \E\{\Vhat(P)\}
  \frac{\chi^2_{\nu_P}}{\nu_P}.
  \label{eq:scaled-chi-working}
\end{equation}
The scale and degrees of freedom are chosen to match the mean and variance
of \(\Vhat(P)\).

For a chi-squared variable with \(\nu_P\) degrees of freedom,
\begin{equation}
  \E(\chi^2_{\nu_P})=\nu_P,
  \qquad
  \Var(\chi^2_{\nu_P})=2\nu_P.
\end{equation}
The scaling in Equation~\eqref{eq:scaled-chi-working} therefore matches the
mean automatically:
\begin{equation}
  \E\left[
  \E\{\Vhat(P)\}\frac{\chi^2_{\nu_P}}{\nu_P}
  \right]
  =\E\{\Vhat(P)\}.
\end{equation}
Its variance is
\begin{align}
  \Var\left[
  \E\{\Vhat(P)\}\frac{\chi^2_{\nu_P}}{\nu_P}
  \right]
  &=
  \frac{[\E\{\Vhat(P)\}]^2}{\nu_P^2}
  \Var(\chi^2_{\nu_P})\\
  &=
  \frac{2[\E\{\Vhat(P)\}]^2}{\nu_P}.
\end{align}
Matching this quantity to \(\Var\{\Vhat(P)\}\) gives
\begin{equation}
  \boxed{
  \nu_P
  =
  \frac{2[\E\{\Vhat(P)\}]^2}
       {\Var\{\Vhat(P)\}}
  }.
  \label{eq:component-df-moments}
\end{equation}
\begin{samepage}
The numerator depends on the mean of the estimated MI variance; the
denominator measures how much that estimate fluctuates between samples.
To first order, \(\E\{\Vhat(P)\}\approx V(P)\), so the remaining task is
to calculate \(\Var\{\Vhat(P)\}\).
\par
\end{samepage}

\begin{samepage}
\section{How the MI Variance Changes with the Table}

To find how much the estimated MI variance changes between samples, we first
ask how the population variance responds when a small amount of probability
is moved into one cell \citep{hampel1974,vanderVaart1998}. Fix a cell
\((x,y)\) and define the perturbed distribution
\begin{equation}
  P_\varepsilon
  =(1-\varepsilon)P+\varepsilon\delta_{(x,y)},
  \label{eq:contamination-path}
\end{equation}
where \(\delta_{(x,y)}\) places probability one on cell \((x,y)\). At
\(\varepsilon=0\), the distribution is \(P\). Increasing \(\varepsilon\)
moves a small amount of probability toward the selected cell while preserving
total probability one.
\par
\end{samepage}

For an arbitrary cell \((i,j)\),
\begin{equation}
  p_{ij}(\varepsilon)
  =(1-\varepsilon)p_{ij}
  +\varepsilon\mathbf 1\{i=x,j=y\}.
\end{equation}
The joint and marginal probability derivatives at the original distribution
are
\begin{align}
  \left.\frac{\mathrm d p_{ij}(\varepsilon)}{\mathrm d\varepsilon}
  \right|_{0}
  &=\mathbf 1\{i=x,j=y\}-p_{ij},
  \label{eq:joint-perturbation-derivative}\\
  \left.\frac{\mathrm d p_{i+}(\varepsilon)}{\mathrm d\varepsilon}
  \right|_{0}
  &=\mathbf 1\{i=x\}-p_{i+},\\
  \left.\frac{\mathrm d p_{+j}(\varepsilon)}{\mathrm d\varepsilon}
  \right|_{0}
  &=\mathbf 1\{j=y\}-p_{+j}.
\end{align}
The derivatives of the joint probabilities sum to zero because
\begin{equation}
  \sum_{i,j}
  \{\mathbf 1(i=x,j=y)-p_{ij}\}
  =1-1=0.
\end{equation}

\medskip
\noindent\textbf{Effect on pointwise MI and MI.}

Along the perturbation path, pointwise MI is
\begin{equation}
  \pmi_{P_\varepsilon}(i,j)
  =\log p_{ij}(\varepsilon)
  -\log p_{i+}(\varepsilon)
  -\log p_{+j}(\varepsilon).
\end{equation}
\begin{samepage}
For each logarithm, the derivative is the probability derivative divided by
the probability. Substituting the joint and marginal derivatives gives
\begin{align}
  \left.
  \frac{\mathrm d}{\mathrm d\varepsilon}
  \pmi_{P_\varepsilon}(i,j)
  \right|_{0}
  &={}
  \frac{\mathbf 1\{i=x,j=y\}-p_{ij}}{p_{ij}}
  -\frac{\mathbf 1\{i=x\}-p_{i+}}{p_{i+}}
  -\frac{\mathbf 1\{j=y\}-p_{+j}}{p_{+j}}
  \notag\\*
  &={}
  \frac{\mathbf 1\{i=x,j=y\}}{p_{ij}}
  -\frac{\mathbf 1\{i=x\}}{p_{i+}}
  -\frac{\mathbf 1\{j=y\}}{p_{+j}}
  +1.
  \label{eq:pmi-derivative}
\end{align}
The joint term supplies \(-1\), while subtracting the row and column
terms supplies \(+1\) each. Their sum is the final \(+1\).
\par
\end{samepage}

MI along the path is
\begin{equation}
  I(P_\varepsilon)
  =\sum_{i,j}p_{ij}(\varepsilon)
  \pmi_{P_\varepsilon}(i,j).
\end{equation}

The product rule separates its derivative into the change in probability
weights and the change in pointwise-MI values:
\begin{align}
  \left.\frac{\mathrm d I(P_\varepsilon)}{\mathrm d\varepsilon}
  \right|_{0}
  ={}&
  \sum_{i,j}
  \{\mathbf 1(i=x,j=y)-p_{ij}\}\pmi_P(i,j)
  \notag\\
  &+
  \sum_{i,j}p_{ij}
  \left.\frac{\mathrm d\pmi_{P_\varepsilon}(i,j)}
  {\mathrm d\varepsilon}\right|_{0}.
\end{align}
The first sum is \(\pmi_P(x,y)-I(P)\). In the second sum, the selected-cell
term contributes \(p_{xy}/p_{xy}=1\). The row term contributes
\(-\sum_j p_{xj}/p_{x+}=-1\); the column term likewise contributes \(-1\);
and the constant contributes \(\sum_{i,j}p_{ij}=1\). These four
contributions cancel. Hence
\begin{equation}
  \boxed{
  \left.\frac{\mathrm d I(P_\varepsilon)}{\mathrm d\varepsilon}
  \right|_{0}
  =\pmi_P(x,y)-I(P)
  }.
  \label{eq:mi-derivative}
\end{equation}
This is the centred pointwise-MI influence function from
Section~\ref{sec:first-order-mi-background}.

\medskip
\noindent\textbf{Effect on the MI variance.}

The target is the first-order change in
\begin{equation}
  V(P)
  =\Var_P\{\pmi_P(X,Y)\}.
\end{equation}
Write this variance in second-moment form:
\begin{equation}
  V(P)=\E_P\{\pmi_P(X,Y)^2\}-I(P)^2.
  \label{eq:v-second-moment}
\end{equation}
Define the observation-level variance sensitivity by
\begin{equation}
  g_P(x,y)
  =
  \left.
  \frac{\mathrm d V(P_\varepsilon)}{\mathrm d\varepsilon}
  \right|_{0}.
  \label{eq:g-definition}
\end{equation}
Differentiating Equation~\eqref{eq:v-second-moment} gives
\begin{equation}
  g_P(x,y)
  =
  \left.\frac{\mathrm d}{\mathrm d\varepsilon}
  \E_{P_\varepsilon}\{\pmi_{P_\varepsilon}(X,Y)^2\}
  \right|_{0}
  -2I(P)
  \left.\frac{\mathrm d I(P_\varepsilon)}{\mathrm d\varepsilon}
  \right|_{0}.
  \label{eq:g-decomposition}
\end{equation}

It remains to differentiate
\begin{equation}
  \E_{P_\varepsilon}\{\pmi_{P_\varepsilon}(X,Y)^2\}
  =\sum_{i,j}p_{ij}(\varepsilon)
  \pmi_{P_\varepsilon}(i,j)^2.
\end{equation}
The product rule gives
\begin{align}
  \left.\frac{\mathrm d}{\mathrm d\varepsilon}
  \E_{P_\varepsilon}\{\pmi_{P_\varepsilon}(X,Y)^2\}
  \right|_{0}
  ={}&
  \sum_{i,j}
  \left.\frac{\mathrm d p_{ij}(\varepsilon)}{\mathrm d\varepsilon}
  \right|_{0}
  \pmi_P(i,j)^2
  \notag\\
  &+
  2\sum_{i,j}p_{ij}\pmi_P(i,j)
  \left.\frac{\mathrm d\pmi_{P_\varepsilon}(i,j)}
  {\mathrm d\varepsilon}\right|_{0}.
  \label{eq:m2-product-rule}
\end{align}
The first sum is
\begin{equation}
  \pmi_P(x,y)^2-\E_P\{\pmi_P(X,Y)^2\}.
  \label{eq:m2-weight-term}
\end{equation}

For the second sum, define the probability-weighted row and column means
\begin{align}
  \E_P\{\pmi_P(X,Y)\mid X=x\}
  &=\frac{\sum_j p_{xj}\pmi_P(x,j)}{p_{x+}},\\
  \E_P\{\pmi_P(X,Y)\mid Y=y\}
  &=\frac{\sum_i p_{iy}\pmi_P(i,y)}{p_{+y}}.
\end{align}
Substitution of Equation~\eqref{eq:pmi-derivative} gives
\begin{align}
  &\sum_{i,j}p_{ij}\pmi_P(i,j)
  \left.\frac{\mathrm d\pmi_{P_\varepsilon}(i,j)}
  {\mathrm d\varepsilon}\right|_{0}
  \notag\\
  &\quad=
  \pmi_P(x,y)
  -\frac{\sum_j p_{xj}\pmi_P(x,j)}{p_{x+}}
  -\frac{\sum_i p_{iy}\pmi_P(i,y)}{p_{+y}}
  +I(P)
  \notag\\
  &\quad=
  \pmi_P(x,y)
  -\E_P\{\pmi_P(X,Y)\mid X=x\}
  -\E_P\{\pmi_P(X,Y)\mid Y=y\}
  +I(P).
  \label{eq:m2-score-term}
\end{align}

Combining Equations~\eqref{eq:mi-derivative},
\eqref{eq:g-decomposition}, \eqref{eq:m2-weight-term}, and
\eqref{eq:m2-score-term} gives the variance sensitivity. To see how its
squared term arises, combine the weight-change term in
Equation~\eqref{eq:m2-weight-term} with the derivative of \(-I(P)^2\):
\begin{align*}
 &\pmi_P(x,y)^2-\E_P\{\pmi_P(X,Y)^2\}
       -2I(P)\{\pmi_P(x,y)-I(P)\}\\*
 &\quad=\pmi_P(x,y)^2-\{V(P)+I(P)^2\}
       -2I(P)\pmi_P(x,y)+2I(P)^2\\*
 &\quad=\pmi_P(x,y)^2-2I(P)\pmi_P(x,y)+I(P)^2-V(P)\\*
 &\quad=\{\pmi_P(x,y)-I(P)\}^2-V(P).
\end{align*}
\begin{samepage}
The second line uses
\(\E_P\{\pmi_P(X,Y)^2\}=V(P)+I(P)^2\); the last line completes the
square. Adding twice the expression in Equation~\eqref{eq:m2-score-term}
gives
\begin{equation}
\boxed{
\begin{aligned}
g_P(x,y)
={}&\{\pmi_P(x,y)-I(P)\}^2-V(P)\\
&+2\bigl[
\pmi_P(x,y)
-\E_P\{\pmi_P(X,Y)\mid X=x\}\\
&\hspace{18mm}
-\E_P\{\pmi_P(X,Y)\mid Y=y\}
+I(P)
\bigr].
\end{aligned}
}
\label{eq:g-final}
\end{equation}
The first line is the change if the pointwise-MI values were held fixed. The
bracket accounts for the fact that they change too: each pointwise-MI value
uses a cell probability and its row and column totals. The conditional means
collect the effects of those changing totals. Both parts are needed to
describe how the estimated variance varies across samples.
\par
\end{samepage}

\section{Component degrees of freedom and final test}

\noindent\textbf{Population component degrees of freedom.}

Define the population variability of the observation-level sensitivities by
\begin{equation}
  \tau^2(P)=\Var_P\{g_P(X,Y)\}.
  \label{eq:tau-definition}
\end{equation}
The sensitivities average to zero under \(P\), as shown in
Appendix~\ref{app:derivation-details}. Their variation, not their average,
determines how much the estimated MI variance changes between samples.

\begin{assumption}[Regular variance sensitivity]
For the first-order Expanded Welch construction,
\(\tau^2(P)>0\) and \(\tau^2(Q)>0\).
\end{assumption}

An observation in cell \(Z_a^{(P)}\) can be represented by a table that places
probability one in that cell, denoted \(\delta_{Z_a^{(P)}}\). Relative to
\(P\), that observation contributes the table change
\(\delta_{Z_a^{(P)}}-P\). The empirical table averages these one-observation
tables, so its difference from \(P\) is
\begin{equation}
  \widehat P-P
  =\frac{1}{n_P}\sum_{a=1}^{n_P}
  \left\{\delta_{Z_a^{(P)}}-P\right\}.
\end{equation}
The plug-in \(\Vhat(P)\) applies the variance calculation to this empirical
table. Equation~\eqref{eq:g-definition} gives its first-order change for
one observation's direction. Because a first-order change is linear in the
table error \citep[Chapter~20]{vanderVaart1998}, the changes from the
observations can be averaged:
\begin{equation}
  \Vhat(P)-V(P)
  \approx
  \frac{1}{n_P}
  \sum_{a=1}^{n_P}g_P(Z_a^{(P)}).
  \label{eq:vhat-linearization}
\end{equation}
The observations are independent, so
\begin{align}
  \Var\{\Vhat(P)\}
  &\approx
  \frac{1}{n_P^2}
  \sum_{a=1}^{n_P}
  \Var_P\{g_P(Z_a^{(P)})\}\\*
  &=\frac{\tau^2(P)}{n_P}.
  \label{eq:vhat-sampling-variance}
\end{align}
Substitution into Equation~\eqref{eq:component-df-moments} gives the
population component degree of freedom
\begin{equation}
  \boxed{
  \nu_P
  \approx
  \frac{2n_PV(P)^2}{\tau^2(P)}
  }.
  \label{eq:population-component-df}
\end{equation}
To interpret Equation~\eqref{eq:population-component-df}, divide the
sampling variance of \(\Vhat(P)\) by the square of the population
quantity it estimates:
\begin{equation}
  \frac{\Var\{\Vhat(P)\}}{V(P)^2}
  \approx \frac{\tau^2(P)}{n_PV(P)^2},
\end{equation}
so the component degrees of freedom are approximately twice the reciprocal of
this relative variance. When the estimate fluctuates little compared with
the size of \(V(P)\), the degrees of freedom are large and the Student
\(t\) distribution is close to the standard normal distribution. Greater
relative fluctuation gives fewer degrees of freedom and a larger rejection
cutoff.

Positive \(V(P)\) does not guarantee positive \(\tau^2(P)\): table changes
can have no first-order effect on \(V(P)\). In that case, describing the
sampling variation of \(\Vhat(P)\) requires a higher-order approximation.

\medskip
\noindent\textbf{Table-based component degrees of freedom.}

The observed table replaces every population quantity in Equation
\eqref{eq:g-final} by its empirical counterpart. The sums below are over
occupied cells, because pointwise MI is undefined at an empirical zero. For
each occupied cell,
\begin{equation}
\begin{aligned}
\widehat g_P(i,j)
={}&\{\widehat\pmi_P(i,j)-\Ihat(P)\}^2-\Vhat(P)\\
&+2\left[
\widehat\pmi_P(i,j)
-\frac{\sum_{j'}\widehat p_{ij'}\widehat\pmi_P(i,j')}
       {\widehat p_{i+}}
-\frac{\sum_{i'}\widehat p_{i'j}\widehat\pmi_P(i',j)}
       {\widehat p_{+j}}
+\Ihat(P)
\right].
\end{aligned}
\label{eq:g-hat}
\end{equation}
In all weighted sums, an empty cell contributes zero; no logarithm or
sensitivity needs to be evaluated there. The empirical variance is centred
explicitly:
\begin{equation}
  \widehat\tau^2(P)
  =\sum_{i,j}\widehat p_{ij}
  \left\{\widehat g_P(i,j)
  -\sum_{k,l}\widehat p_{kl}\widehat g_P(k,l)\right\}^2.
  \label{eq:tau-hat}
\end{equation}
The estimated component degree of freedom is
\begin{equation}
  \boxed{
  \widehat\nu_P
  =\frac{2n_P\Vhat(P)^2}{\widehat\tau^2(P)}
  }.
  \label{eq:nu-p-hat}
\end{equation}
Repeating the same calculation for the second table gives
\(\widehat\nu_Q\).

\medskip
\noindent\textbf{Combining the two components.}

Substituting the two variance contributions and component degrees of freedom
into Equation~\eqref{eq:ws-general} gives
\begin{equation}
  \boxed{
  \widehat\nu_{\mathrm{expanded}}
  =
  \frac{
  \left\{\Vhat(P)/n_P+\Vhat(Q)/n_Q\right\}^2
  }{
  \left\{\Vhat(P)/n_P\right\}^2/\widehat\nu_P
  +
  \left\{\Vhat(Q)/n_Q\right\}^2/\widehat\nu_Q
  }
  }.
  \label{eq:expanded-df}
\end{equation}
The expanded two-sided p-value is
\begin{equation}
  p_{\mathrm{expanded}}
  =2\left\{1-F_{t_{\widehat\nu_{\mathrm{expanded}}}}(|T|)\right\},
  \label{eq:expanded-p}
\end{equation}
where $F_{t_\nu}$ denotes the cumulative distribution function of a Student
$t$ variable with $\nu$ degrees of freedom.

For a fixed finite \(T\), the Student \(t\) distribution gives a two-sided
p-value at least as large as the standard normal distribution used by Wald.
Under the implemented validity rules, every valid Expanded Welch result is
also Wald-valid. Thus every
Expanded Welch rejection is also a Wald rejection, including when invalid
outputs count as non-rejections in the reported rates. The correction cannot
raise rejection where Wald is already too conservative.
Section~\ref{sec:zero-component} examines a zero-component case in the
combined degrees-of-freedom formula.

\section{Worked example and implementation}

Consider the two observed \(3\times3\) tables
\begin{equation}
  N^P=
  \begin{pmatrix}
  35&10&5\\
  10&20&5\\
  5&5&10
  \end{pmatrix},
  \qquad
  N^Q=
  \begin{pmatrix}
  30&5&5\\
  5&30&5\\
  5&5&10
  \end{pmatrix}.
  \label{eq:worked-tables}
\end{equation}
Their totals are \(n_P=105\) and \(n_Q=100\). In the first table, cell
\((1,1)\) has count 35, and its row and column each total 50. Its empirical
pointwise MI is therefore
\begin{equation*}
 \widehat\pmi_P(1,1)
 =\log\left\{\frac{35/105}{(50/105)(50/105)}\right\}
 =\log(1.47)\simeq0.38526.
\end{equation*}
This cell contributes \((35/105)(0.38526)\) to the plug-in MI estimate.
Summing the contributions from all nine cells gives
\(\Ihat(P)=0.13693\); the same calculation for the second table gives
\(\Ihat(Q)=0.25848\) nats. Since \((3-1)(3-1)=4\), the corrected difference is
\begin{equation}
  \Deltahat
  =\left(0.13693-\frac{4}{2(105)}\right)
  -\left(0.25848-\frac{4}{2(100)}\right)
  =-0.12059.
\end{equation}

The variance calculation weights each squared deviation from the estimated
MI by its cell probability. Cell \((1,1)\), for example, contributes
\((35/105)(0.38526-0.13693)^2\simeq0.02056\) to \(\Vhat(P)\).
Summing over all cells gives \(\Vhat(P)=0.27189\) and
\(\Vhat(Q)=0.43371\). Hence
\begin{equation}
  \SEhat
  =\left(\frac{0.27189}{105}+\frac{0.43371}{100}\right)^{1/2}
  =0.08323,
  \qquad
  T=-1.44898.
\end{equation}
Normal Wald gives \(p=0.14734\).

For Expanded Welch, the same cell illustrates the sensitivity calculation.
The first row's pointwise-MI values are 0.38526, \(-0.51083\) and
\(-0.64436\), weighted by its counts 35, 10 and 5. Its conditional mean is
\begin{equation*}
 \frac{35(0.38526)+10(-0.51083)+5(-0.64436)}{50}
 \simeq0.10308.
\end{equation*}
The first column has the same conditional mean because this count table is
symmetric. Substituting these values into Equation~\eqref{eq:g-hat} gives
\begin{align*}
 \widehat g_P(1,1)
 &\simeq (0.38526-0.13693)^2-0.27189\\*
 &\quad+2\{0.38526-0.10308-0.10308+0.13693\}\\*
 &\simeq0.4218.
\end{align*}
The probability-weighted mean of the sensitivities is zero, so this cell
contributes \((35/105)(0.4218)^2\simeq0.0593\) to
\(\widehat\tau^2(P)\). Repeating the calculation for all cells in both
tables gives \(\widehat\tau^2(P)=0.93212\) and
\(\widehat\tau^2(Q)=0.82035\). Substitution into
Equation~\eqref{eq:nu-p-hat} gives component degrees of freedom
\begin{equation}
  \widehat\nu_P=16.65,
  \qquad
  \widehat\nu_Q=45.86.
\end{equation}
Combining them with Equation~\eqref{eq:expanded-df} gives
\(\widehat\nu_{\mathrm{expanded}}=59.03\) and \(p=0.15264\).
Neither method rejects at the 5\% level.

\pagebreak[3]
\medskip
\noindent\textbf{Operational algorithm and complexity.}
\label{sec:expanded-complexity}

Table~\ref{tab:expanded-algorithm} summarises one comparison.

\begin{table}[htbp]
\centering
\caption{Operational Expanded Welch--Satterthwaite algorithm.}
\label{tab:expanded-algorithm}
\begin{tabularx}{\textwidth}{@{}cX@{}}
\toprule
Step & Calculation \\
\midrule
1 & Validate the two aligned count tables and their declared alphabet dimensions; independence of samples is a design assumption. \\
2 & Compute empirical probabilities, pointwise MI, plug-in MI, and the leading bias-corrected difference. \\
3 & Compute \(\Vhat(P)\), \(\Vhat(Q)\), the shared standard error, and \(T\). \\
4 & Compute the row and column pointwise-MI means needed in \(\widehat g_P\) and \(\widehat g_Q\). \\
5 & Compute \(\widehat\tau^2(P)\), \(\widehat\tau^2(Q)\), and the two component degrees of freedom. \\
6 & Combine the components using Equation~\eqref{eq:expanded-df}. \\
7 & Return the p-value from the Student \(t\) distribution, effective degrees of freedom, and validity diagnostics. \\
\bottomrule
\end{tabularx}
\end{table}

Each empirical probability, pointwise-MI value, row mean, column mean, and
variance contribution is calculated by a fixed number of table scans. The
time complexity is therefore \(O(rc)\), and the working memory is also
\(O(rc)\).

\section{Behaviour and limitations}
\label{sec:near-independence-mechanism}

\noindent\textbf{Near independence.}

Consider a population table that
approaches independence while its row and column margins stay fixed. At
independence, MI is zero and its first-order change is also zero. The
leading changes in MI, the pointwise-MI variance \(V(P)\), and the variance
sensitivity \(\tau^2(P)\) can therefore be compared using a second-order
expansion. Appendix~\ref{sec:fixed-margin-local-derivation} gives the
calculation and its fixed-margin assumptions. Close to independence on
that path, it gives
\begin{equation}
 V(P)\approx 2I(P),
 \qquad \tau^2(P)\approx4V(P).
 \label{eq:near-independent-moments}
\end{equation}
Substituting these local approximations into
the component degrees-of-freedom formula gives
\begin{equation}
 \nu_P=\frac{2n_PV(P)^2}{\tau^2(P)}
 \approx\frac{n_PV(P)}{2}
 \approx n_PI(P).
 \label{eq:near-independent-df}
\end{equation}
For a fixed sample size, this approximation makes the degrees of freedom
smaller as population MI approaches zero. The cutoff from the Student
\(t\) distribution then increases, so a larger absolute statistic is
needed to reject equal MI.

These approximations require the fixed-margin path, fixed table size and
positive cell probabilities. Using them for samples also requires a useful
first-order sampling approximation. At exact independence, quadratic sampling
variation replaces the first-order term; Section~\ref{sec:independence-boundary}
derives that case.

\medskip
\noindent\textbf{Theoretical status and validity conditions.}

Equations~\eqref{eq:mi-derivative} and \eqref{eq:g-final} are exact
derivatives for positive population tables. They provide first-order
approximations to the sampling errors in MI and its estimated variance.

The scaled chi-squared model in Equation~\eqref{eq:scaled-chi-working} and
the Student \(t\) distribution for \(T\) are finite-sample approximations.
The MI difference and standard error come from the same tables. A sample
that changes the estimated difference can also change its standard error,
and the degrees-of-freedom adjustment does not account for how the two move
together.
Appendix~\ref{sec:mi-variance-covariance} gives the first-order covariance
between each population's MI and variance estimates.

As sample sizes grow under regular conditions, the estimated
degrees of freedom increase and the Student \(t\) distribution approaches the
standard normal distribution used by Wald.

Both tests need a finite corrected MI difference and a positive, finite
estimated standard error. Expanded Welch also needs positive, finite component
and combined degrees of freedom. A failed requirement gives an invalid
calculation, not evidence for the null.
Appendix~\ref{sec:validity-rules} specifies the numerical criteria;
the experiment evaluates validity alongside calibration.
